\subsubsection{Adjacent-order quotient dynamics and an obstruction to scalar closure}
\label{subsec:adjacentratio}
The preterminal spatial condition
\eqref{eq:openslabrh} would imply RH, but constructing its
coefficient by dividing by $U_n$ presupposes precisely
the positivity that remains to be established.
There is a useful exact PDE for these quotients
\emph{where positivity has independently been proved}.
It also explains why a tempting order-monotone
bootstrap cannot close this gap.

Put $k_n=(n+1)(2n+1)$,
$D_n=U_n$, $V_n=k_nU_{n+1}$,
$W_n=k_{n+1}U_{n+2}$ and
$\mathcal D=\partial_s-\frac12\partial_x^2$.

\begin{theorem}[Division-free adjacent-order identity]
\label{thm:adjacentratio}
For the genuine theta heat hierarchy,
\begin{equation}
\boxed{D_n\mathcal D V_n-V_n\mathcal D D_n
=V_n^2-k_nD_nW_n.}
\label{eq:adjacentdivisionfree}
\end{equation}
On every open region on which $D_n,V_n>0$,
set $R_n=V_n/D_n$ and
$R_{n+1}=k_nW_n/V_n$.
Then
\begin{equation}
\boxed{
\partial_s R_n=\frac12\partial_x^2R_n
+(\partial_x\log U_n)\partial_xR_n
+R_n(R_n-R_{n+1}).}
\label{eq:adjacentheat}
\end{equation}
\end{theorem}
\begin{proof}
The hierarchy gives $\mathcal D D_n=-V_n$
and $\mathcal D V_n=-k_nW_n$.
Their antisymmetric product is
\eqref{eq:adjacentdivisionfree}.
Substitute $V_n=R_nD_n$ into that identity
and use the product rule
$\mathcal D(R_nD_n)=D_n\mathcal DR_n+
R_n\mathcal DD_n-(\partial_xR_n)(\partial_xD_n)$.
Division by the positive $D_n^2$ yields
\eqref{eq:adjacentheat}.
\end{proof}

If the as-yet-unproved ordering
$R_n\le R_{n+1}$ held, the reaction term
would be nonpositive, suggesting a
scalar comparison. That ordering is
false for general real-zero transforms.

\begin{proposition}[Exact adjacent-order reversal]
\label{prop:ratioreversal}
For the polynomial $G(z)=z(z-1)(z-10)$,
the generalized Laguerre ratios at the simple
zero $z=0$ satisfy
\[
\boxed{R_1(0)=\frac{303}{50}>
\frac{15}{101}=R_2(0).}
\]
\end{proposition}
\begin{proof}
At a simple real root $a$, let
$b_j=(a-\rho_j)^{-2}$ for the other
real roots, and let $e_k$ be the
elementary symmetric polynomials in
the $b_j$. The product
$G(a+iy)G(a-iy)=G'(a)^2y^2
\prod_j(1+b_jy^2)$ gives
$R_n(a)=k_ne_n/e_{n-1}$.
Here $b_1=1$, $b_2=1/100$, so
$e_1=101/100$ and $e_2=1/100$,
which gives the displayed exact fractions.
\end{proof}

As a diagnostic on the \emph{genuine} theta
heat family at $t=1/2$, direct
$46$-decimal-digit Fourier quadrature gives
\begin{center}
\begin{tabular}{rcc}
\hline
$x$ & $R_1$ & $R_2$\\
\hline
$0$ & $0.0116731531$ & $0.0337807219$\\
$20$ & $0.0277960654$ & $0.0559941838$\\
$28.269$ & $12.005931723$ & $0.1041071217$\\
$42.044$ & $7.126195805$ & $0.2152913414$\\
\hline
\end{tabular}
\end{center}
The values suggest a similar reversal near
known heat-flow real zeros; they are not
directed interval certificates and are
not used to establish a theta inequality.

For a Gaussian Fourier transform with
$A(s)=3+2s$, the exact ratio
$R_n(s,x)=(2n+1)/A(s)$ satisfies
\eqref{eq:adjacentheat} with a negative reaction.
For the quadratic collision control
$H_t(z)=z^2-2t$, by contrast,
$R_1=3/(x^2+1-2s)$ and $R_2=0$;
the reaction is positive and increases
towards the terminal collision.
Both models pass exact rational
or symbolic checks.

The next genuine analytic target is therefore
to bound the transport
$\partial_x\log U_n$ and the signed reaction
\emph{without} assuming positivity of
the denominators. A proposed upper barrier
$R_n\lesssim1+\log^2(2+|x|)$ is not closed
until those bounds and spatial boundary
conditions are established. The division-free
identity \eqref{eq:adjacentdivisionfree}
remains well defined at zeros and may be a
better starting point; alternatively one must
prove the theta-specific Duhamel or dyadic
defect bound already isolated above.
This is a sharpened proof obstruction,
not a proof of RH.

The exact algebra appears in
\path{XiKernel/AdjacentLaguerreRatio.lean}.
The symbolic and numerical reproduction
is in \path{certificates/xi_adjacent_ratio_heat_exact.py}.
The Lean statement does not formalize the
complex Fourier integrals or full
parabolic quotient calculation.
